\documentclass[a4paper,12pt]{article}
\usepackage[T2A]{fontenc}
\usepackage[utf8]{inputenc}
\usepackage[russian]{babel}
\usepackage{amsmath,amssymb,mathtools}
\usepackage{helvet}
\renewcommand{\familydefault}{\sfdefault}
\usepackage{icomma}
\usepackage[a4paper,left=21mm,right=21mm,top=23mm,bottom=23mm,headheight=23pt]{geometry}
\usepackage{microtype}
\usepackage{tikz}
\usetikzlibrary{calc,arrows.meta,positioning,shapes.geometric}
\usepackage{booktabs,tabularx,array}
\usepackage{enumitem}
\usepackage{hyperref}
\usepackage{parskip}
\usepackage{fancyhdr}
\usepackage{setspace}
\usepackage{xcolor}
\definecolor{accent}{HTML}{008080}
\definecolor{darkaccent}{HTML}{004D4D}
\definecolor{textgray}{HTML}{333333}
\definecolor{softgray}{HTML}{747474}
\hypersetup{colorlinks=true,linkcolor=darkaccent,urlcolor=darkaccent}
\setstretch{1.11}
\setlength{\parindent}{0pt}
\setlength{\parskip}{3pt}
\setlist[itemize]{leftmargin=18pt,itemsep=2pt,topsep=3pt}
\setlist[enumerate]{leftmargin=21pt,itemsep=2pt,topsep=2pt}
\renewcommand{\arraystretch}{1.25}
\pagestyle{fancy}
\fancyhf{}
\fancyhead[L]{\footnotesize\color{softgray}Связанные тела и второй закон Ньютона}
\fancyhead[R]{\footnotesize\color{softgray}Володя Хачатурян}
\fancyfoot[C]{\footnotesize\thepage}
\renewcommand{\headrulewidth}{0.35pt}
\renewcommand{\headrule}{\hbox to\headwidth{\color{accent}\leaders\hrule height \headrulewidth\hfill}}
\newcommand{\mys}[1]{\section*{\textcolor{darkaccent}{#1}}\vspace{-7pt}\noindent{\color{accent}\rule{\linewidth}{0.4pt}}\par\vspace{2pt}}
\newcommand{\subt}[1]{\par\vspace{6pt}{\color{darkaccent}\bfseries #1}\par\vspace{-1pt}}
\newcommand{\checkitem}[1]{\item $\square$\enspace #1}
\newcommand{\nterm}{\,\text{Н}}
\newcommand{\aunit}{\,\text{м/с}^{2}}
\begin{document}
\begin{titlepage}
\thispagestyle{empty}
\centering
\vspace*{22mm}
{\color{darkaccent}\large ФИЗИКА\par}
\vspace{16mm}
{\LARGE Чек-лист\par}
\vspace{11mm}
{\color{accent}\LARGE Связанные тела\par}
\vspace{2mm}
{\color{accent}\Large и второй закон Ньютона\par}
\vspace{18mm}
{\normalsize Силы, проекции, нить и блоки\par}
\vfill
{\large Володя Хачатурян\par}
\vspace{14mm}
{19 ноября 2026 г.\par}
\vspace{22mm}
{\small Лёвин Артём Александрович\par}
{\small эксклюзивно для Володи Хачатуряна\par}
\vspace{18mm}
\begin{tikzpicture}
\draw[accent,thin] (0,0) .. controls (1.45,.55) and (2.6,-.45) .. (4.0,0)
 .. controls (5.3,.45) and (6.6,-.5) .. (8.0,0);
\end{tikzpicture}
\vspace{3mm}
\end{titlepage}
\setcounter{page}{1}

\mys{1. С чего начинается решение}
\textcolor{darkaccent}{\textbf{Главная идея.}} Рассматривайте каждое тело отдельно: нарисуйте силы, выберите положительные направления осей, запишите второй закон Ньютона в проекциях, затем свяжите уравнения условиями о нити.

\subt{Направления четырёх основных сил}
\begin{tabularx}{\linewidth}{@{}p{37mm}p{33mm}X@{}}
\toprule
\textbf{Сила} & \textbf{Обозначение} & \textbf{Куда направлена}\tabularnewline
\midrule
Тяжести & $\vec F_{\text{тяж}}=m\vec g$ & Строго вертикально вниз.\tabularnewline
Реакции опоры & $\vec N$ & Перпендикулярно поверхности опоры, от неё.\tabularnewline
Трения скольжения & $\vec F_{\text{тр}}$ & Против относительного движения соприкасающихся поверхностей.\tabularnewline
Натяжения нити & $\vec T$ & Вдоль натянутой нити, от тела к нити.\tabularnewline
\bottomrule
\end{tabularx}
\textcolor{darkaccent}{Важно:} на наклонной плоскости сила реакции опоры обычно направлена под углом к вертикали, поэтому её нельзя автоматически считать «вертикально вверх».

\subt{Второй закон Ньютона}
\[
\sum\vec F=m\vec a,\qquad
\sum F_x=ma_x,\qquad \sum F_y=ma_y.
\]
Знак каждой проекции определяется направлением \emph{выбранной оси}: вдоль положительного направления --- плюс, против --- минус. Если ускорение направлено против оси, в правой части пишите $-ma$.

\subt{Свойства идеальной нити и блока}
\begin{itemize}
\item \textbf{Невесомая нить:} при идеальном (без трения, с пренебрежимо малой массой) блоке натяжение одинаково на участках одной нити: $T_1=T_2=T$.
\item \textbf{Нерастяжимая натянутая нить:} связанные тела в рассмотренных схемах имеют равные \emph{модули} ускорения: $|a_1|=|a_2|=a$. Направления и знаки проекций могут различаться.
\item \textbf{Трение скольжения:} $F_{\text{тр}}=\mu N$, где $\mu$ --- коэффициент трения. Если в условии сказано «без трения», $F_{\text{тр}}=0$.
\end{itemize}

\mys{2. Два тела на горизонтальной поверхности}
\begin{center}
\begin{tikzpicture}[>=Latex,scale=.90, every node/.style={font=\small}]
\draw[textgray,line width=.8pt] (.2,.3)--(12,.3);
\foreach \x in {.35,.85,...,11.85} {\draw[softgray,thin] (\x,.3)--++(-.24,-.22);}
\draw[fill=gray!9,draw=textgray,thick] (2,.3) rectangle (3.6,1.48);
\draw[fill=gray!9,draw=textgray,thick] (6.5,.3) rectangle (8.1,1.48);
\node at (2.8,.86) {$m_2$};\node at (7.3,.86) {$m_1$};
\draw[thick] (3.6,.92)--(6.5,.92);
\draw[->,very thick,accent] (8.1,1.05)--(10.0,1.05) node[right] {$\vec F$};
\draw[->,thick,darkaccent] (7.3,1.7)--(8.7,1.7) node[right] {$\vec a$};
\draw[->,thick,darkaccent] (2.8,1.7)--(4.2,1.7) node[right] {$\vec a$};
\node[below] at (5.05,.89) {нить};
\end{tikzpicture}
\end{center}
Тело 1 тянут вправо силой $F$; тело 2 движется вслед за ним. Ось $x$ направим вправо, ось $y$ --- вверх. Если оба тела скользят вправо, силы трения направлены влево.
\[
\begin{aligned}
\text{Тело 1:}\quad &F-T-F_{\text{тр}1}=m_1a, & N_1-m_1g&=0,\\
\text{Тело 2:}\quad &T-F_{\text{тр}2}=m_2a, & N_2-m_2g&=0.
\end{aligned}
\]
Учитываем $F_{\text{тр}1}=\mu N_1=\mu m_1g$ и $F_{\text{тр}2}=\mu m_2g$ (при одинаковом $\mu$). Складываем горизонтальные уравнения: натяжение сокращается.
\[
F-\mu(m_1+m_2)g=(m_1+m_2)a,
\qquad a=\frac{F-\mu(m_1+m_2)g}{m_1+m_2}.
\]
\textcolor{darkaccent}{\textbf{Пример.}} $m_1=4$ кг, $m_2=2$ кг, $F=24\nterm$, $\mu=0{,}1$, $g=10\aunit$.
\[
a=\frac{24-0{,}1\cdot6\cdot10}{6}=3\aunit,
\qquad T=m_2a+\mu m_2g=8\nterm.
\]

\mys{3. Два груза на неподвижном блоке}
\begin{center}
\begin{tikzpicture}[>=Latex,scale=.90,every node/.style={font=\small}]
\draw[softgray,thick] (4,3.2)--(4,2.65);
\draw[fill=gray!12,thick] (4,2.25) circle (.43);
\fill[white] (4,2.25) circle (.07);
\draw[very thick] (3.57,.93)--(3.57,2.25) arc[start angle=180,end angle=0,radius=.43]--(4.43,.62);
\draw[fill=gray!9,thick] (3.1,.38) rectangle (4.04,.93);
\draw[fill=gray!9,thick] (3.96,.07) rectangle (4.9,.62);
\node at (3.57,.65) {$m_1$}; \node at (4.43,.35) {$m_2$};
\draw[->,thick,accent] (2.7,.25)--(2.7,1.35) node[above] {$\vec a_1$};
\draw[->,thick,accent] (5.3,1.25)--(5.3,.15) node[below] {$\vec a_2$};
\node[align=left,anchor=west] at (5.9,2.3) {Блок идеальный,\\ось неподвижна};
\end{tikzpicture}
\end{center}
Пусть $m_2>m_1$. Груз 1 ускоряется вверх, груз 2 --- вниз. Для \emph{обоих грузов} направим ось $y$ вверх. Тогда знаки ускорений различны:
\[
\underbrace{T-m_1g=m_1a}_{\text{груз 1}},
\qquad
\underbrace{T-m_2g=-m_2a}_{\text{груз 2}}.
\]
Вычитая второе уравнение из первого, получаем
\[
a=\frac{(m_2-m_1)g}{m_1+m_2},\qquad
T=m_1(g+a)=m_2(g-a).
\]
Если окажется $a<0$, фактическое движение противоположно принятому предположению.

\clearpage
\mys{4. Наклонная плоскость и подвешенный груз}
\begin{center}
\begin{tikzpicture}[>=Latex,scale=.88,every node/.style={font=\small}]
\coordinate (A) at (0,0);\coordinate (B) at (5.3,0);\coordinate (C) at (5.3,3.06);
\draw[draw=textgray,thick] (A)--(C)--(B)--cycle;
\draw[fill=gray!12,draw=textgray,thick,rotate around={30:(2.95,1.7)}] (2.36,1.43) rectangle (3.54,1.97);
\node[rotate=30,fill=white,inner sep=0.5pt] at (2.95,1.7) {$m_1$};
\draw[fill=gray!12,thick] (5.62,2.78) circle (.24);
\draw[thick] (3.63,2.09)--(5.48,3.0) arc[start angle=150,end angle=0,radius=.24]--(5.86,.83);
\draw[fill=gray!9,thick] (5.45,.1) rectangle (6.27,.83);
\node at (5.86,.47) {$m_2$};
\draw[->,very thick,accent] (1.52,1.14)--(2.2,1.53) node[above] {$\vec a_1$};
\draw[->,very thick,accent] (6.65,2.05)--(6.65,.95) node[below] {$\vec a_2$};
\draw[accent] (.8,0) arc[start angle=0,end angle=30,radius=.8];
\node at (1.05,.22) {$\alpha$};
\end{tikzpicture}
\end{center}
Для первого тела выбираем $x$ \emph{вдоль плоскости вверх}, $y$ --- \emph{перпендикулярно плоскости от неё}; для висящего тела используем вертикальную ось вверх. Предположим, что $m_1$ движется вверх по плоскости, $m_2$ --- вниз.
\[
\begin{aligned}
\text{Тело 1, ось }x:&&T-m_1g\sin\alpha-F_{\text{тр}}&=m_1a,\\
\text{Тело 1, ось }y:&&N-m_1g\cos\alpha&=0,\\
\text{Тело 2, ось }y:&&T-m_2g&=-m_2a.
\end{aligned}
\]
Здесь $N=m_1g\cos\alpha$, а при скольжении вверх $F_{\text{тр}}=\mu m_1g\cos\alpha$. Если трения нет, полагаем $F_{\text{тр}}=0$. Вычитая уравнение для висящего груза из уравнения вдоль плоскости, получаем
\[
 a=\frac{m_2g-m_1g\sin\alpha-F_{\text{тр}}}{m_1+m_2}.
\]
\textcolor{darkaccent}{\textbf{Пример без трения.}} $m_1=0{,}3$ кг, $m_2=0{,}2$ кг, $\alpha=30^\circ$, $g=10\aunit$.
\[
\begin{aligned}
 T-0{,}3\cdot10\cdot\sin30^\circ&=0{,}3a, &T-1{,}5&=0{,}3a,\\
 T-0{,}2\cdot10&=-0{,}2a, &T-2&=-0{,}2a.
\end{aligned}
\]
Вычитаем второе уравнение из первого: $0{,}5=0{,}5a$, поэтому
\[
\boxed{a=1\aunit},\qquad
\boxed{T=0{,}3\cdot1+1{,}5=1{,}8\nterm}.
\]
Положительное значение $a$ подтверждает выбранное направление движения.

\mys{5. Как избегать типичных ошибок}
\begin{enumerate}
\item На наклонной плоскости вдоль склона действует $mg\sin\alpha$, перпендикулярно склону --- $mg\cos\alpha$.
\item Ось направлена вверх, груз ускоряется вниз: в правой части уравнения стоит $-ma$.
\item $N=mg$ на горизонтальной опоре без других вертикальных сил; на наклонной --- $N=mg\cos\alpha$.
\item Натяжение направлено вдоль нити от тела; у идеальной нити одинаковы \emph{модули} натяжений.
\item Контроль: $[a]=\text{м/с}^2$, $[T]=\text{Н}$; $a<0$ означает обратное направление ускорения.
\end{enumerate}

\clearpage
\thispagestyle{plain}
\begin{center}
{\LARGE\color{darkaccent}Алгоритм решения задачи}\par
\vspace{3mm}
{\normalsize Связанные тела, нити и блоки}\par
\vspace{13mm}
\begin{tikzpicture}[
 >=Latex,
 action/.style={rectangle,rounded corners=2mm,draw=accent,thick,align=center,minimum width=93mm,minimum height=12mm,text width=85mm,inner sep=3.5mm,font=\small},
 smallaction/.style={rectangle,rounded corners=2mm,draw=accent,thick,align=center,minimum width=42mm,minimum height=11mm,text width=36mm,inner sep=2mm,font=\small},
 decision/.style={diamond,draw=darkaccent,thick,aspect=3.3,align=center,inner sep=1.5mm,text width=35mm,font=\small},
 terminal/.style={ellipse,draw=darkaccent,thick,align=center,minimum width=53mm,minimum height=11mm,font=\small},
 arrow/.style={->,thick,color=textgray}]
\node[terminal] (start) at (0,0) {Начало};
\node[action] (a) at (0,-1.75) {Изобразите каждое тело отдельно и укажите все действующие силы};
\node[action] (b) at (0,-3.65) {Выберите оси; определите проекции сил и ускорений со знаками};
\node[decision] (d) at (0,-6.02) {Есть трение скольжения?};
\node[smallaction] (yes) at (4.48,-7.60) {$F_{\text{тр}}=\mu N$};
\node[smallaction] (no) at (-4.48,-7.60) {$F_{\text{тр}}=0$};
\node[action] (eq) at (0,-9.10) {Для каждого тела запишите $\sum F_x=ma_x$ и $\sum F_y=ma_y$};
\node[action] (thread) at (0,-10.95) {Примените условия идеальной нити: одинаковое $T$ и равные модули $a$};
\node[action] (solve) at (0,-12.80) {Решите систему уравнений; найдите ускорение и натяжение};
\node[decision] (check) at (0,-15.13) {Знаки и размерности верны?};
\node[smallaction] (fix) at (-4.48,-16.80) {Исправьте проекции или знаки};
\node[terminal] (end) at (0,-17.55) {Ответ};
\draw[arrow] (start)--(a);
\draw[arrow] (a)--(b);
\draw[arrow] (b)--(d);
\draw[arrow] (d.east)-| node[pos=.25,above] {да} (yes.north);
\draw[arrow] (d.west)-| node[pos=.25,above] {нет} (no.north);
\draw[arrow] (yes.south)|- (eq.east);
\draw[arrow] (no.south)|- (eq.west);
\draw[arrow] (eq)--(thread);
\draw[arrow] (thread)--(solve);
\draw[arrow] (solve)--(check);
\draw[arrow] (check.south)--node[right] {да} (end.north);
\draw[arrow] (check.west)-| node[pos=.25,above] {нет} (fix.north);
\draw[arrow] (fix.west)--++(-.5,0)|- (b.west);
\end{tikzpicture}
\end{center}
\clearpage
\mys{Тренировочный блок}
\textcolor{darkaccent}{\textbf{Путь А --- 10 заданий.}} В задачах считайте $g=10\aunit$. Нити невесомы и нерастяжимы, блоки идеальны, если иное не оговорено. Обосновывайте знаки проекций.

\subt{Базовый уровень}
\begin{enumerate}
\item Брусок скользит вправо по горизонтальному столу. К нему справа прикреплена натянутая горизонтальная нить. Укажите направления сил тяжести, реакции опоры, трения скольжения и натяжения нити.
\item Брусок массой $2$ кг скользит по горизонтальной поверхности с коэффициентом трения $\mu=0{,}2$. Определите модуль реакции опоры и силу трения скольжения.
\item Два бруска массами $m_1=3$ кг и $m_2=2$ кг связаны нитью на гладком горизонтальном столе. К первому бруску приложена сила $F=15\nterm$ вправо. Найдите ускорение и натяжение нити.
\end{enumerate}
\subt{Средний уровень}
\begin{enumerate}[start=4]
\item Два бруска $m_1=3$ кг и $m_2=2$ кг движутся вправо по горизонтали; коэффициент трения у обоих $\mu=0{,}1$. На первое тело вправо действует $F=25\nterm$. Найдите ускорение и натяжение нити.
\item Через неподвижный идеальный блок перекинута нить с грузами $m_1=1$ кг и $m_2=3$ кг. Найдите ускорение и натяжение нити. Укажите направление ускорения каждого груза.
\item В системе с идеальным неподвижным блоком груз $m_1=2$ кг ускоряется вверх с модулем $2\aunit$. Найдите массу второго груза и натяжение нити.
\item На гладкой наклонной плоскости с углом $30^\circ$ находится брусок $m_1=2$ кг, связанный через идеальный блок с висящим грузом $m_2=1$ кг. Определите ускорение системы и натяжение нити.
\end{enumerate}
\subt{Повышенный уровень}
\begin{enumerate}[start=8]
\item Брусок $m_1=0{,}6$ кг находится на гладкой плоскости под углом $30^\circ$ и связан через блок с грузом $m_2=0{,}4$ кг. Определите модуль и направление ускорения, а также натяжение нити.
\item Брусок $m_1=2$ кг лежит на наклонной плоскости. Известно, что $\sin\alpha=0{,}6$, $\cos\alpha=0{,}8$, коэффициент трения скольжения $\mu=0{,}25$. Через блок он соединён с грузом $m_2=2$ кг. Найдите ускорение и натяжение нити, если второй груз движется вниз.
\item Брусок $m_1=4$ кг движется вверх по гладкой наклонной плоскости под углом $30^\circ$, а подвешенный груз $m_2$ движется вниз с ускорением $2\aunit$. Определите $m_2$ и натяжение нити.
\end{enumerate}

\clearpage
\mys{Ответы к тренировочному блоку}
\begin{tabularx}{\linewidth}{@{}p{12mm}X@{}}
\toprule
\textbf{№} & \textbf{Ответ}\\
\midrule
1 & Сила тяжести --- вниз; реакция опоры --- вверх; трение --- влево; натяжение --- вправо.\\
2 & $N=20\nterm$; $F_{\text{тр}}=4\nterm$.\\
3 & $a=3\aunit$; $T=6\nterm$.\\
4 & $a=4\aunit$; $T=10\nterm$.\\
5 & $a=5\aunit$; $T=15\nterm$. Груз $m_2$ --- вниз, $m_1$ --- вверх.\\
6 & $m_2=3$ кг; $T=24\nterm$.\\
7 & $a=0$; $T=10\nterm$ (система в равновесии, если первоначально покоилась).\\
8 & $a=1\aunit$; $T=3{,}6\nterm$; $m_2$ --- вниз, $m_1$ --- вверх по плоскости.\\
9 & $a=1\aunit$; $T=18\nterm$; второй груз --- вниз.\\
10 & $m_2=3{,}5$ кг; $T=28\nterm$.\\
\bottomrule
\end{tabularx}
\par\vspace{9mm}
\textcolor{darkaccent}{\textbf{Итоговая памятка}}\par
\begin{tabularx}{\linewidth}{@{}p{51mm}X@{}}
\toprule
\textbf{Сигнал в условии} & \textbf{Ваше действие}\\
\midrule
«Нить невесомая» & При идеальном блоке приравняйте модули натяжений.\\
«Нить нерастяжимая» & Приравняйте модули ускорений; проверьте их знаки в осях.\\
«Гладкая плоскость» & Положите силу трения равной нулю.\\
«Наклонная плоскость» & Выберите оси вдоль плоскости и перпендикулярно ей.\\
«Найти ускорение» & Запишите уравнение для каждого тела и решите систему.\\
«Получилось $a<0$» & Пересмотрите предполагаемое направление ускорения.\\
\bottomrule
\end{tabularx}
\vfill
\begin{center}
\small\color{softgray}Лёвин Артём Александрович · эксклюзивно для Володи Хачатуряна
\end{center}
\end{document}
